\begin{initiativkarte}{AG Antifa}

  % Bild (optional)
  \IfFileExists{bilder/ag_antifa.jpg}{%
    \begin{center}
      \includegraphics[width=\linewidth,height=5cm,keepaspectratio]{bilder/ag_antifa.jpg}
    \end{center}
    \vspace{0.4cm}
  }{}

  % Kurzbeschreibung
  \textit{\textcolor{hawgray}{Die AG Antifa beschäftigt sich mit der Aufarbeitung der Hochschulgeschichte sowie mit antifaschistischer Bildungs- und Erinnerungsarbeit.}}

  \vspace{0.5cm}
  \hawrule

  % Beschreibung
  \textbf{\textcolor{hawblue}{Was machen wir?}}\\[0.3cm]

  Die AG Antifa der HAW Hamburg setzt sich nach eigener Darstellung mit der historischen Rolle von Hochschulen im Nationalsozialismus auseinander und beschäftigt sich mit antifaschistischer Bildungs- und Erinnerungsarbeit.

  Im Mittelpunkt stehen dabei politische und historische Bildung sowie der kritische Blick auf gesellschaftliche Strukturen und deren Entwicklung.

  Zu den Aktivitäten gehören:
  \begin{itemize}
      \item politische und historische Bildungsarbeit
      \item Diskussionen zu gesellschaftlichen Themen
      \item Vernetzung mit engagierten Studierenden
      \item Auseinandersetzung mit Hochschul- und Zeitgeschichte
  \end{itemize}

  Die Gruppe bietet einen Raum für Austausch, Reflexion und kritische Diskussionen über Geschichte und Gegenwart.

  \vspace{0.5cm}

  % Tags
  \tag{Politik}
  \tag{Geschichte}
  \tag{Bildung}
  \tag{Diskussion}
  \tag{Gesellschaft}

  \vspace{0.6cm}
  \hawrule

  % Infozeilen
  \infozeile{\faUsers}{Zielgruppe: Studierende aller Fachrichtungen}
  \infozeile{\faGraduationCap}{Semester: Ab dem 1. Semester}
  \infozeile{\faEnvelope}{Kontakt: \href{mailto:antifa-haw@outlook.com}{antifa-haw@outlook.com}}
  \infozeile{\faGlobe}{Instagram: \href{https://www.instagram.com/antifa_hawhamburg}{@antifa\_hawhamburg}}

\end{initiativkarte}